\documentclass[UTF8, 11pt]{ctexart}

\usepackage[margin=1in]{geometry}
\usepackage{amsmath, amssymb}
\usepackage{booktabs}
\usepackage{graphicx}
\usepackage{hyperref}
\usepackage{url}

\title{\textbf{实验报告}\\
面向通信高效分布式优化的可扩展图划分算法}
\author{小组成员：姓名 1，姓名 2，姓名 3}
\date{CS240：算法设计与分析，2026 年春季学期}

\begin{document}

\maketitle

\section{背景与相关工作}

本项目研究图划分算法在通信高效分布式优化中的应用。在分布式优化中，大规模
目标函数通常由多个智能体或 工作节点协同求解，并通过通信网络交换信息。如果
强相关的变量或智能体被分配到不同 工作节点上，跨工作节点 通信开销就会增加。
因此，一个好的划分应当在保持计算负载均衡的同时，尽量减少跨分区依赖边。

图划分是经典算法问题，在并行科学计算、稀疏矩阵计算、分布式图处理和大规模
机器学习中都有应用。本项目关注平衡 $k$ 路图划分问题。相关方法包括谱划分、
局部搜索细化和多层图划分。主要复现方法是 Kernighan--Lin 启发式算法，它通过
移动或交换顶点迭代降低割值。我们将其与谱二分方法和贪心均衡基线进行比较。

\section{代表性方法分析}

设 $G=(V,E,w)$ 是一个无向加权图。顶点表示数据块、变量块、智能体或计算任务，
边表示依赖关系或通信链路。给定分区数 $k$，目标是找到互不相交且并集为 $V$ 的
子集 $V_1,\ldots,V_k$：
\[
    \min_{V_1,\ldots,V_k}
    \sum_{(u,v)\in E,\;u\in V_i,\;v\in V_j,\;i\neq j} w_{uv},
\]
并满足负载均衡约束
\[
    |V_i| \leq (1+\epsilon)\frac{|V|}{k}, \quad i=1,\ldots,k.
\]
目标函数表示跨工作节点 通信代价，均衡约束表示工作负载分配。

贪心基线逐个分配顶点，在满足容量约束的分区中选择局部割代价较低的分区。谱
划分构造图拉普拉斯矩阵 $L=D-A$，利用 Fiedler 向量进行二分，并递归得到 $k$ 个
分区。Kernighan--Lin 细化从初始划分出发，对分区对执行局部搜索。谱方法需要
特征分解；在本项目实现中，稠密矩阵特征分解复杂度为 $O(n^3)$，对课程规模实验
是可接受的。贪心基线约为 $O(mk)$，细化复杂度取决于迭代轮数和成对局部子问题。

\section{复现与实现}

代码以 Python 包的形式组织在 \path{src/} 目录下，主要模块包括：
\begin{itemize}
    \item \path{src/graphs.py}：合成图和小型真实图生成；
    \item \path{src/algorithms.py}：贪心、谱划分和 Kernighan--Lin 方法；
    \item \path{src/metrics.py}：割边权重、归一化割和负载均衡比例；
    \item \path{src/simulation.py}：一致性平均通信仿真；
    \item \path{src/experiments.py}：实验入口和 CSV 结果导出；
    \item \path{src/plotting.py}：实验图像生成。
\end{itemize}

复现的核心部分是 Kernighan--Lin 细化思想。我们使用 NetworkX 中稳定的成对
Kernighan--Lin 二分过程，并将其嵌入模块化的 $k$ 路细化循环中，初始划分由谱
划分给出。这是对 $k$ 路平衡图划分的一种实践化改写。所有实验均使用固定随机
种子，便于复现。

\section{实验评估}

实验使用网格图、Erdos--Renyi 随机图、随机几何图和 Zachary karate club 图。
评价指标包括割边权重、负载均衡比例、运行时间，以及一致性仿真中的通信量。
原始结果保存在 \path{outputs/partition_results.csv}。

在全部 benchmark 上，平均割边权重为：
\[
\text{greedy}=98.25,\quad
\text{spectral}=72.75,\quad
\text{Kernighan--Lin}=66.88.
\]
一致性仿真中的平均跨分区消息数为：
\[
\text{greedy}=15460,\quad
\text{spectral}=11360,\quad
\text{Kernighan--Lin}=10440.
\]
因此，Kernighan--Lin 细化方法取得了最低的平均割值和最低的平均跨工作节点 通信
量。在八个图实例中，它有六个实例取得最佳割值。谱划分在网格图和随机几何图
上表现也很好，因为 Fiedler 向量能够捕捉这类图的空间结构。

\begin{figure}[ht]
    \centering
    \includegraphics[width=0.75\linewidth]{outputs/figures/cut_weight.png}
    \caption{割边权重比较。数值越低，表示跨工作节点 通信越少。}
\end{figure}

\begin{figure}[ht]
    \centering
    \includegraphics[width=0.75\linewidth]{outputs/figures/runtime_seconds.png}
    \caption{运行时间比较。贪心基线最快，Kernighan--Lin 需要额外细化时间。}
\end{figure}

\begin{figure}[ht]
    \centering
    \includegraphics[width=0.75\linewidth]{outputs/figures/balance_ratio.png}
    \caption{负载均衡比例比较。测试实例中，各方法都保持了较好的均衡性。}
\end{figure}

同一张图上，不同划分的一致性误差几乎相同。这是因为仿真保持原始通信图不变，
划分只用于统计跨工作节点 消息数。因此，划分改变的是通信放置方式，而不是底层
平均动力学。这一结果符合预期：图划分在该应用中的主要价值体现在减少跨分区
通信开销。

\section{可选拓展与讨论}

当前项目在纯图划分指标之外，加入了一个轻量级的分布式一致性仿真。主要限制是
该仿真没有改变通信拓扑本身，而只是统计消息是否跨工作节点 边界。进一步拓展可以
加入通信预算、跨工作节点 延迟或加权通信代价。另一个自然拓展是在更大图上与
METIS 等完整多层图划分器进行比较。

总体而言，实验支持本项目的核心观点：经典图划分算法可以在现代分布式优化场景
中用于降低通信成本，同时保持较好的负载均衡。

\end{document}

